% Generated by Listas/make_statements.py from pratica-lecture-08.md, item 12. Do not edit.
\textbf{Source.} Tadelis (2013), \emph{Game Theory: An Introduction}, Exercise 16.4 ``School Choice'', pp.~336--337 {[}PDF pp.~355--356{]}. Reproduced verbatim. Answer: none in the book.



\begin{quote}

\textbf{School Choice:} A prospective student, player 1, is deciding whether to go to Brown Bear University (\(B\)) or Silly Tree University (\(S\)). Both are top-rated universities that accept only the very best students. The difference is that Silly Tree offers a more holistic and comfortable lifestyle, while Brown Bear requires harder work and imposes more ``studying'' costs on its students. For this reason a player who goes to Brown Bear learns to become more self-sufficient and becomes more productive in the workforce. The cost of learning and level of final productivity depend on the type of player 1, who can be either excellent (\(E\)) or just really good (\(G\)) (because no other type would be accepted into these universities). Player 1 knows his type, but all others in society know only that a proportion \(p\) of young adults are type \(E\). The cost of learning and the level of productivity from each choice are given as follows:



{\def\LTcaptype{none} % do not increment counter

\begin{longtable}[]{@{}cccc@{}}

\toprule\noalign{}

Type & University choice & Learning cost & Productivity \\

\midrule\noalign{}

\endhead

\bottomrule\noalign{}

\endlastfoot

\(E\) & \(B\) & 2 & 12 \\

\(E\) & \(S\) & 0 & 4 \\

\(G\) & \(B\) & 8 & 10 \\

\(G\) & \(S\) & 2 & 2 \\

\end{longtable}

}



Once player 1 finishes school, he is hired by a firm (player 2), which can place him in one of two jobs: low-tech (\(L\)) or high-tech (\(H\)). The wage for the \(L\) job is \(w_L = 2\) and that for the \(H\) job is \(w_H = 6\). The payoffs to player 1 are the wages less the cost of education. The firm's profits depend on both the job assignment and the type of employee. If the employee is assigned to an \(H\) job, the net profits to the firm are equal to the productivity of the employed player 1 less the wage he is paid. If the assignment is to an \(L\) job, the net profits are half the productivity of the employed player 1 less the wage he is paid.



\begin{enumerate}

\def\labelenumi{\alph{enumi}.}

\item

  Draw this game in extensive form.

\item

  Assume that \(p = \frac{1}{2}\). Represent the matrix form of the Bayesian game.

\item

  Find all the pure-strategy Bayesian Nash equilibria.

\item

  Find all the pure-strategy perfect Bayesian equilibria.

\item

  In brief terms, what is the intuition that explains the comparison between your results in (d) and (c)?

\end{enumerate}

\end{quote}

